\section{Weryfikacja hipotez}
\label{sec:hipotezy}

\subsection{H1: przewaga Valkey i MongoDB rośnie z liczbą klientów}

Dla każdego rozmiaru i~każdej pary (baza nierelacyjna, baza relacyjna)
liczymy przewagę jako iloraz przepustowości przy strategiach domyślnych, dla
1, 10, 100 i~1000 klientów. Przewaga ,,rośnie'', jeśli współczynnik korelacji
rang Spearmana między liczbą klientów a~przewagą jest dodatni i~przewaga przy
1000 klientów jest większa niż przy jednym.

\begin{table}[H]
\centering
\caption{H1: przewaga przepustowości (iloraz) przy 1 i~przy 1000 klientach.}
\label{tab:h1}
\small
\tbinput{generated/tables/h1.tex}
\end{table}

Przewaga rośnie w~\HOneGrows{} z~\HOneCases{} przypadków. Wynik
automatycznej oceny: hipoteza H1 \textbf{\HOneVerdict}.

Iloraz przepustowości rośnie z~liczbą klientów dla wszystkich par i~rozmiarów.
Współczynnik Spearmana wynosi wszędzie 1 (zależność ściśle monotoniczna). Dla
10 mln biletów przewaga Valkey nad PostgreSQL rośnie
z~\cratio{valkey}{postgres}{10000000}{lua}{cond}{1}$\times$ przy jednym kliencie
do \cratio{valkey}{postgres}{10000000}{lua}{cond}{1000}$\times$ przy tysiącu.
Nad MariaDB rośnie z~\cratio{valkey}{mariadb}{10000000}{lua}{cond}{1}$\times$
do \cratio{valkey}{mariadb}{10000000}{lua}{cond}{1000}$\times$.

Wynik dla MongoDB trzeba czytać ostrożniej. Iloraz MongoDB do PostgreSQL rośnie
z~\cratio{mongo}{postgres}{10000000}{atomic}{cond}{1} do
\cratio{mongo}{postgres}{10000000}{atomic}{cond}{1000}: przy małej
współbieżności MongoDB jest kilka razy wolniejszy, a~przy 1000 klientach
dorównuje PostgreSQL. Względna pozycja MongoDB rośnie, ale realną przewagę
ma on tylko nad MariaDB. Hipotezę H1 uznajemy za
\textbf{potwierdzoną dla Valkey}. Dla MongoDB potwierdza się teza
o~rosnącym stosunku wydajności, a~nie o~przewadze w~każdym punkcie.

Dwa czynniki wzmacniają wynik Valkey i~trzeba je brać pod uwagę: brak zapisu
na dysk oraz łączenie poleceń w~potoki przez sterownik. Strategia \texttt{watch}
pokazuje też, że ten sam silnik z~nieodpowiednią metodą synchronizacji
przegrywa ze wszystkimi bazami relacyjnymi. H1 jest więc prawdziwa dla
właściwie zaprojektowanej operacji atomowej, a~nie dla silnika
niezależnie od sposobu użycia.


\subsection{H2: różnice między silnikami rosną z liczbą rekordów}

Dla każdego scenariusza i~wariantu indeksów liczymy iloraz czasu
najwolniejszej i~najszybszej bazy. Różnica ,,rośnie'', jeśli ten iloraz przy
10 mln biletów jest większy niż przy 500 tys.

\begin{table}[H]
\centering
\caption{H2: iloraz czasu najwolniejszej i~najszybszej bazy.}
\label{tab:h2}
\footnotesize
\IfFileExists{generated/tables/h2.tex}{\begin{minipage}{\linewidth}\centering% generated by tbench report
\begin{tabular}{llrrrr}
\toprule
Scen. & indeksy & 500 tys. & 1 mln & 10 mln & rośnie? \\
\midrule
C1 & bez & \num{4.89} & \num{4.42} & \num{9.17} & tak \\
C1 & z & \num{5.57} & \num{13.0} & \num{12.5} & tak \\
C2 & bez & \num{7.38} & \num{3.13} & \num{7.09} & nie \\
C2 & z & \num{6.17} & \num{6.70} & \num{6.84} & tak \\
C3 & bez & \num{5.55} & \num{4.59} & \num{6.01} & tak \\
C3 & z & \num{8.86} & \num{6.81} & \num{9.00} & tak \\
C4 & bez & \num{2.98} & \num{3.83} & \num{4.17} & tak \\
C4 & z & \num{4.58} & \num{4.16} & \num{4.43} & nie \\
C5 & bez & \num{3.06} & \num{1.57} & \num{1.58} & nie \\
C5 & z & \num{3.70} & \num{2.37} & \num{3.43} & nie \\
C6 & bez & \num{8.92} & \num{5.43} & \num{8.02} & nie \\
C6 & z & \num{13.3} & \num{9.17} & \num{12.0} & nie \\
R1 & bez & \num{447.8} & \num{776.1} & \num{24351} & tak \\
R1 & z & \num{1.45} & \num{1.79} & \num{12.0} & tak \\
R2 & bez & \num{81.5} & \num{160.0} & \num{1954} & tak \\
R2 & z & \num{1.60} & \num{2.20} & \num{1.82} & tak \\
R3 & bez & \num{3.56} & \num{4.51} & \num{15.6} & tak \\
R3 & z & \num{2.52} & \num{2.20} & \num{2.60} & tak \\
R4 & bez & \num{137.6} & \num{73.4} & \num{83.4} & nie \\
R4 & z & \num{1.92} & \num{2.08} & \num{1.60} & nie \\
R5 & bez & \num{511.8} & \num{551.6} & \num{1036} & tak \\
R5 & z & \num{80.2} & \num{80.9} & \num{258.0} & tak \\
R6 & bez & \num{826.0} & \num{1434} & \num{15256} & tak \\
R6 & z & \num{923.9} & \num{1784} & \num{17619} & tak \\
U1 & bez & \num{7.37} & \num{10.7} & \num{4.82} & nie \\
U1 & z & \num{5.37} & \num{5.25} & \num{12.4} & tak \\
U2 & bez & \num{7.21} & \num{13.1} & \num{11.2} & tak \\
U2 & z & \num{8.02} & \num{6.20} & \num{5.33} & nie \\
U3 & bez & \num{67.3} & \num{119.6} & \num{1069} & tak \\
U3 & z & \num{3.19} & \num{2.68} & \num{1.77} & nie \\
U4 & bez & \num{5.17} & \num{4.60} & \num{5.16} & nie \\
U4 & z & \num{3.45} & \num{3.45} & \num{4.65} & tak \\
U5 & bez & \num{3.31} & \num{3.98} & \num{5.49} & tak \\
U5 & z & \num{3.48} & \num{7.40} & \num{5.18} & tak \\
U6 & bez & \num{16.7} & \num{20.1} & \num{20.4} & tak \\
U6 & z & \num{3.04} & \num{2.46} & \num{2.86} & nie \\
D1 & bez & \num{858.5} & \num{239.6} & \num{16378} & tak \\
D1 & z & \num{7.82} & \num{7.09} & \num{9.76} & tak \\
D2 & bez & \num{771.7} & \num{416.6} & \num{16771} & tak \\
D2 & z & \num{8.11} & \num{6.46} & \num{9.42} & tak \\
D3 & bez & \num{240.8} & \num{264.0} & \num{358.0} & tak \\
D3 & z & \num{23.0} & \num{17.8} & \num{17.6} & nie \\
D4 & bez & \num{23.3} & \num{18.8} & \num{20.7} & nie \\
D4 & z & \num{5.05} & \num{6.53} & \num{60.3} & tak \\
D5 & bez & \num{1170} & \num{2152} & \num{26468} & tak \\
D5 & z & \num{94.4} & \num{107.8} & \num{110.8} & tak \\
D6 & bez & \num{20.2} & \num{22.3} & \num{17.0} & nie \\
D6 & z & \num{11.0} & \num{11.4} & \num{9.61} & nie \\
J1 & bez & \num{3.76} & \num{4.24} & \num{8.11} & tak \\
J1 & z & \num{5.51} & \num{3.02} & \num{7.00} & tak \\
J2 & bez & \num{235.6} & \num{275.3} & \num{417.1} & tak \\
J2 & z & \num{16.1} & \num{23.2} & \num{28.3} & tak \\
J3 & bez & \num{5.40} & \num{7.78} & \num{8.19} & tak \\
J3 & z & \num{11.0} & \num{10.4} & \num{9.98} & nie \\
\bottomrule
\end{tabular}
\end{minipage}}{brak danych}
\end{table}

Różnica rośnie w~\HTwoGrows{} z~\HTwoCases{} przypadków. Wynik automatycznej
oceny: hipoteza H2 \textbf{\HTwoVerdict}.

Rozrzut między najszybszą a~najwolniejszą bazą rośnie z~rozmiarem w~większości
przypadków, ale nie we wszystkich. Rośnie tam, gdzie co najmniej jedna baza
skanuje dane w~czasie $O(n)$, a~inna ma ścieżkę stałą lub logarytmiczną.
Przykłady to scenariusze bez indeksów R1, R2, R5, U3, D1, D2, D5 i~J2 oraz
R6, gdzie Valkey odczytuje licznik, a~pozostałe bazy agregują wszystkie bilety
(tab.~\ref{tab:h2}). Przy R4, D4 i~D6 bez indeksów rozrzut nie rośnie, bo
wszystkie bazy, także najszybsza, skanują tam pełne dane. Czas najszybszej
rośnie więc co najmniej w~tym samym tempie co czas najwolniejszej (Valkey ze
skanem \texttt{SCAN}).
Na przykład w~R1 bez indeksów MariaDB jest
\ratio{mariadb}{valkey}{500000}{off}{R1}$\times$ wolniejsza od Valkey przy
500 tys. biletów i~\ratio{mariadb}{valkey}{10000000}{off}{R1}$\times$ przy
10 mln.

Rozrzut nie rośnie dla operacji punktowych po kluczu (C1, C4, U1, U2, J3),
których czas w~każdej bazie jest stały. Nie rośnie też dla zapytań z~indeksem
o~małej selektywności, gdzie wszystkie bazy czytają kilka wierszy. H2 jest
\textbf{potwierdzona częściowo}. Różnice rosną z~liczbą rekordów wtedy, gdy
bazy mają różną złożoność ścieżki dostępu, czyli przy braku indeksu, przy
różnicy modeli danych (osadzenie a~złączenie) lub przy agregacji
utrzymywanej w~locie. Dla dobrze zaindeksowanych operacji punktowych rozmiar
zbioru nie ma znaczenia.

